\documentclass[10pt,a4paper]{amsart}
\usepackage[margin=19mm]{geometry}
\usepackage{amsmath,amssymb,mathtools}
\usepackage[scr=rsfs,cal=euler]{mathalfa}
\usepackage{xfrac}
\usepackage[unicode,hidelinks,bookmarks=false]{hyperref}
\newcommand{\ie}{i.e.\ }
\newcommand{\cf}{cf.\ }
\newcommand{\resp}{respectively\ }
\newcommand{\Subsection}{\subsection}
\providecommand{\nobreakdash}{\nobreak\hbox{-}}
\setlength{\emergencystretch}{2em}
\multlinegap0pt
\usepackage{amssymb}
\usepackage{enumerate}
\newtheorem{thm}{Theorem}
\title{\bf A Reduction of the Reconstruction Conjecture using Domination and Vertex Pair Parameters}
\author{J. Antony Aravind, S. Monikandan}
\date{Source: arXiv:2601.00620v1 (CC BY 4.0)}
\begin{document}
\maketitle

\noindent\emph{Source and licence.} \bf A Reduction of the Reconstruction Conjecture using Domination and Vertex Pair Parameters. Version arXiv:2601.00620v1 (CC BY 4.0), distributed by arXiv under the Creative Commons Attribution 4.0 International licence, http://creativecommons.org/licenses/by/4.0/, as recorded in the arXiv licence metadata for this exact version and shown on the abstract page https://arxiv.org/abs/2601.00620v1. Full licence notice: \texttt{licenses/arxiv-2601.00620-CC-BY-4.0.txt}.\par\smallskip

\noindent\textit{Editorial scope.} Counted unit: the article's thm t3 (source0.tex lines 280--286). The article's complete original proof of it is delivered with it (source0.tex lines 287--324, one \texttt{proof} environment of 38 lines). No mathematics is written, repaired or re-derived: the statement and the proof are the article's own lines, in the article's own notation, and no retained line is substituted or re-flowed.  No further source-local context is needed: every cross-reference of every retained line resolves inside the two delivered spans.  Nothing was omitted from any retained span, so the package carries no omission record.  No retained line contains a citation, so the wrapper carries no bibliography and no bibliography entry.  No retained line contains a figure environment, an image-inclusion command or drawing code, so no image is delivered and the package declares no asset dependency.  Editorial formatting only, in addition: an AMS \texttt{amsart} preamble that restates the article's own declarations and macros used by the retained lines, namely the environments \texttt{thm} on separate counters, together with the standard packages those lines need.  The retained environments print in this document as Theorem 1 in order of appearance, whereas the article prints the same items with the numbering of its own full document.  The complete native source of the article as distributed in the arXiv e-print for this exact version is bundled unchanged as \texttt{sources/192009/source0.tex}.
\begin{thm}\label{t3}
For a connected graph $G$ on $n$ vertices and $k_2 \neq k_3,$ 
\begin{multline} 
\sum_{i=1}^{n}{dv(G-v_i,k_1,k_2,k_3)}=(n-2-k_1-k_2-k_3)dv(G,k_1,k_2,k_3)\\+(k_1+1)dv(G,k_1+1,k_2,k_3)
+(k_2+1)dv(G,k_1,k_2+1,k_3)+(k_3+1)dv(G,k_1,k_2,k_3+1). 
\end{multline}
\end{thm}

\begin{proof}
To prove this, we evaluate the left-hand side, $\sum_{i=1}^{n}{dv(G-v_i,k_1,k_2,k_3)},$ by counting how pairs of non-adjacent vertices with specific neighbour properties in $G$ contribute to $dv(G-v_i,k_1,k_2,k_3)$ across all cards $G-v_i.$ Recall that $dv(G,k_1,k_2,k_3)$ denotes the number of non-adjacent vertex pairs $\{u,v\}$ in $G$  with $k_1$ vertices adjacent to both $u$ and $v,$ $k_2$ vertices adjacent only to $u,$ and $k_3$ vertices adjacent only to $v.$ We consider four types of pairs in $G$ and their contributions when a vertex is deleted:
\begin{enumerate}
\item[(1)] Pairs with $(k_1,k_2,k_3)$ in $G$:
 Let $\{u,v\}$ be a pair of non-adjacent vertices in $G$ where $k_1$ vertices adjacent to both $u$ and $v,$ $k_2$ vertices adjacent only to $u,$ and $k_3$ vertices adjacent only to $v.$ In the deck $\mathscr{D},$ consider the cards:

\begin{enumerate}
\item[{\bf \LARGE .}] $G-u$ and $G-v:$ The pair  $\{u,v\}$ is absent (since one vertex is deleted).
\item[{\bf \LARGE .}] Cards $G-w$ where $w$ is one of the $k_1+k_2+k_3$  vertices adjacent to $u$ or $v:$ Deleting $w$ reduces the count of neighbours, so the pair $\{u,v\}$ no longer has exactly $k_1,k_2,k_3$ neighbours in $G-w$, and thus does not contribute to $dv(G-v_i,k_1,k_2,k_3)$. 
\item[{\bf \LARGE .}] Remaining cards: There are $n$ vertices total, and we exclude $u$, $v,$ and the $k_1+k_2+k_3$ neighbours, leaving $n-(2+k_1+k_2+k_3)$ cards ( noting $u$ and $v$ are distinct, and $k_2 \neq k_3$ ensures asymmetry). In each of these, $\{u,v\}$ remains a non-adjacent pair with $k_1,k_2,k_3$ unchanged, contributing 1 to $dv(G-v_i,k_1,k_2,k_3)$. Thus, each such pair contributes to $n-(2+k_1+k_2+k_3)$ cards, and the total contribution is $(n-2-k_1-k_2-k_3)dv(G,k_1,k_2,k_3).$
\end{enumerate}
\item[(2)] Pairs with $(k_1+1,k_2,k_3)$ in $G$:\\
Consider a pair $\{u,v\}$ in $G$ with $k_1+1$ common neighbours, $k_2$ vertices adjacent only to $u,$ and $k_3$ vertices adjacent only to $v.$ In $G-w,$ where $w$ is one of the $k_1+1$ common neighbours:
\begin{enumerate}
\item[{\bf \LARGE .}] Deleting $w$ reduces the common neighbours to $k_1,$ while $k_2$ and $k_3$ remain unchanged, so $\{u,v\}$ becomes a $(k_1,k_2,k_3)$-pair in $G-w,$ contributing 1 to $dv(G-w,k_1,k_2,k_3).$
\item[{\bf \LARGE .}] In the remaining $n-(k_1+1)$ cards, the pair either disappears (if $u$ or $v$ is deleted) or has $k_1+1$ common neighbours, not $k_1.$ There are $k_1+1$ such cards, so the total contribution is $(k_1+1)dv(G, k_1+1,k_2,k_3).$
\end{enumerate}

\item[(3)] Pairs with $(k_1,k_2+1,k_3)$ in $G$:\\
Take a pair $\{u,v\}$ with $k_1$ common neighbours, $k_2+1$ vertices adjacent only to $u,$ and $k_3$ vertices adjacent only to $v,$  consider $G-w$ where $w$ is one of the $k_3+1$ vertices adjacent only to $v:$
\begin{enumerate}
\item[{\bf \LARGE .}] Deleting $w$ reduces the count to $k_2$ vertices adjacent only to $u,$ while $k_1$ and $k_3$ stay the same, making  $\{u,v\}$ a $(k_1,k_2,k_3)$-pair in $G-w.$
\item[{\bf \LARGE .}] In other cards, the pair either vanishes or retains $k_2+1.$ With $k_2+1$ such cards, the contribution is $(k_2+1)dv(G, k_1,k_2+1,k_3).$
\end{enumerate}

\item[(4)] Pairs with $(k_1,k_2,k_3+1)$ in $G$:\\
For a pair $\{u,v\}$ with $k_1$ common neighbours, $k_2$ vertices adjacent only to $u,$ and $k_3+1$ vertices adjacent only to $v:$
\begin{enumerate}
\item[{\bf \LARGE .}] Deleting $w$ adjusts the count to $k_3,$ with $k_1$ and $k_2$ unchanged, so $\{u,v\}$ fits $(k_1,k_2,k_3)$ in $G-w.$
\item[{\bf \LARGE .}] Elsewhere, the pair either disappears or keeps $k_3+1.$ This occurs in $k_3+1$ cards, yielding $(k_3+1)dv(G, k_1,k_2,k_3+1).$
\end{enumerate}
\end{enumerate}
Summing these contributions, no other pairs in $G$ produce $k_1,k_2,k_3$ in $G-v_i$ under single vertex deletion (since $k_2 \neq k_3$ ensures distinct roles), so:
\begin{multline} \sum_{i=1}^{n}{dv(G-v_i,k_1,k_2,k_3)}=(n-2-k_1-k_2-k_3)dv(G,k_1,k_2,k_3)\\+(k_1+1)dv(G,k_1+1,k_2,k_3)
+(k_2+1)dv(G,k_1,k_2+1,k_3)+(k_3+1)dv(G,k_1,k_2,k_3+1). 
\end{multline}

\end{proof}

\end{document}
